% TOP_PROBLEM: {"id": 30002215, "title": "$\\mathcal Z$-Stability and Regularity of Nuclear $C^*$-Algebras", "category_id": 4, "category": {"id": 4, "name": "algebra", "display_name": "Algebra", "description": "Group theory, ring theory, field theory, and algebraic structures.", "slug": "algebra", "order_index": 4, "created_at": "2026-07-31T15:26:25.671Z"}, "status": "open"}
% FIRST_POSTED: 2026-10-06
% ATTEMPT_STATUS: unresolved
% OPENAI_CLAIM: {"schema_version":1,"source_repo":"https://github.com/openai/math","source_commit":"adc7f1241b42e322a6451854ab7e4b4c146bf78a","release_date":"2026-10-06","match":"partial","resolution":"partial","independently_reviewed":false,"second_pass":"n/a","families":[{"family":"291","title":"Cuntz comparison, nuclear dimension, and equivariant Jiang–Su stability","subject":"Operator algebras","manuscripts":[{"title":"Nuclear dimension and Jiang–Su stability without elementary subquotients","pdf_path":"preprints/Nuclear-dimension-and-Jiang-Su-stability-without-elementary-subquotients-September-23-2026/paper.pdf","readme_path":"preprints/Nuclear-dimension-and-Jiang-Su-stability-without-elementary-subquotients-September-23-2026/README.md","date":"2026-09-23","theorem_ref":"Introduction main regularity theorems"},{"title":"Tracial projection methods and uniform property Gamma","pdf_path":"preprints/Tracial-projection-methods-and-uniform-property-Gamma-September-23-2026/paper.pdf","readme_path":"preprints/Tracial-projection-methods-and-uniform-property-Gamma-September-23-2026/README.md","date":"2026-09-23","theorem_ref":"Introduction main regularity theorems"}],"lean":{"doc_path":"lean/docs/291.md","comparators":["lean/ComparatorChallenges/UniformGamma.json","lean/ComparatorChallenges/UniformGamma.lean"],"declarations":["OAI.ComparatorModel.CurrentMain.main"],"covers_main_theorem":false},"reasoning_trace":null}]}
\documentclass[11pt]{article}
\usepackage[margin=25mm]{geometry}
\usepackage{fontspec}
\setmainfont{Latin Modern Roman}
\usepackage{amsmath,amssymb,mathtools}
\usepackage{unicode-math}
\setmathfont{Latin Modern Math}
\usepackage{enumitem}
\usepackage{xurl,hyperref,bookmark}
\hypersetup{colorlinks=true,urlcolor=blue}
\setcounter{secnumdepth}{0}
\setlength{\parindent}{0pt}
\setlength{\parskip}{5pt}
\setlength{\emergencystretch}{3em}
\begin{document}
\section*{\$\textbackslash{}mathcal Z\$-Stability and Regularity of Nuclear \$C\textasciicircum{}*\$-Algebras}

\subsection{OpenAI's claim}
Manuscript 1 (family 291, 2026-09-23): Finite nuclear dimension is claimed to characterize Jiang--Su stability for separable nuclear algebras lacking elementary ideal subquotients; moreover, Jiang--Su stabilization always lowers nuclear dimension to at most one. Manuscript 2 (family 291, 2026-09-23): The paper claims uniform Gamma from real rank zero in the tracial-completion ultrapower for simple unital nuclear stably finite algebras. Comparison gives absorption, also under specified bounded-trace hypotheses in the stably projectionless setting.

\subsection{Outline of the argument}
Manuscript 1 (family 291, 2026-09-23): Localized order-zero approximations control ranks even when traces are infinite on some ideals. A supporting map and comparison yield two-color approximations that survive stabilization and hereditary passage. Conversely, Gaussian sign tensors and pure-state excision construct two full orthogonal central elements from the absence of characters, supplying the absorption criterion.

Manuscript 2 (family 291, 2026-09-23): Nuclear averaging and scalar block compressions preserve coefficient traces while packing small projections with controlled commutators. A separate comparison construction rounds ranks, forms clipped sums of centered symmetries, and extracts Gaussian trace partitions. In the projectionless case, a combined overlap-and-model cost controls small splittings before reindexing amplifies them to exact central partitions.

\subsection{Scope relative to this problem}
Manuscript 1 (family 291, 2026-09-23): The paper resolves regularity for the specified separable nonelementary or infinite-dimensional setting, but the broader strict-comparison question in this bundle supplies no such complete hypothesis restriction.

Manuscript 2 (family 291, 2026-09-23): The paper resolves regularity for the specified separable nonelementary or infinite-dimensional setting, but the broader strict-comparison question in this bundle supplies no such complete hypothesis restriction.

\subsection{Formal verification}
For family 291, the scope page documents the following selections. The selection gives uniform Gamma from real rank zero of the designated tracial ultrapower. The documented comparison scope also treats Jiang--Su absorption and Gamma under strict comparison, with additional trace and target-rank conditions in the projectionless case. Manuscript 1 is not among the papers selected by that page, so the family link does not formalize its main conclusion. The documented statement covers the main claim of manuscript 2. Comparator files: UniformGamma.lean. Declarations: OAI.ComparatorModel.CurrentMain.main.

This site has not executed the Lean files or checked their dependencies. These descriptions report the scope of the linked declarations, rather than a new verification of the proofs.

\subsection{Status caveat}
OpenAI attributes these results to an internal model. Its release notice cautions that results lacking full formalization can contain mistakes. This site has not independently assessed the proofs; a solved label records the scope of the reported claim.

\subsection{Sources}
\begin{itemize}
\item OpenAI, \emph{Nuclear dimension and Jiang–Su stability without elementary subquotients}, 2026-09-23; citation key \texttt{OAI:Nuclear-dimension-and-Jiang-Su-stability-without-elementary-subquotients-September-23-2026}. \href{https://github.com/openai/math/blob/adc7f1241b42e322a6451854ab7e4b4c146bf78a/preprints/Nuclear-dimension-and-Jiang-Su-stability-without-elementary-subquotients-September-23-2026/paper.pdf}{PDF (GitHub)}; \href{https://github.com/openai/math/raw/adc7f1241b42e322a6451854ab7e4b4c146bf78a/preprints/Nuclear-dimension-and-Jiang-Su-stability-without-elementary-subquotients-September-23-2026/paper.pdf}{PDF (direct)}; \href{https://github.com/openai/math/blob/adc7f1241b42e322a6451854ab7e4b4c146bf78a/preprints/Nuclear-dimension-and-Jiang-Su-stability-without-elementary-subquotients-September-23-2026/README.md}{Manuscript page and citation}; \href{https://github.com/openai/math/blob/main/preprints/Nuclear-dimension-and-Jiang-Su-stability-without-elementary-subquotients-September-23-2026/paper.pdf}{Latest version}.
\item OpenAI, \emph{Tracial projection methods and uniform property Gamma}, 2026-09-23; citation key \texttt{OAI:Tracial-projection-methods-and-uniform-property-Gamma-September-23-2026}. \href{https://github.com/openai/math/blob/adc7f1241b42e322a6451854ab7e4b4c146bf78a/preprints/Tracial-projection-methods-and-uniform-property-Gamma-September-23-2026/paper.pdf}{PDF (GitHub)}; \href{https://github.com/openai/math/raw/adc7f1241b42e322a6451854ab7e4b4c146bf78a/preprints/Tracial-projection-methods-and-uniform-property-Gamma-September-23-2026/paper.pdf}{PDF (direct)}; \href{https://github.com/openai/math/blob/adc7f1241b42e322a6451854ab7e4b4c146bf78a/preprints/Tracial-projection-methods-and-uniform-property-Gamma-September-23-2026/README.md}{Manuscript page and citation}; \href{https://github.com/openai/math/blob/main/preprints/Tracial-projection-methods-and-uniform-property-Gamma-September-23-2026/paper.pdf}{Latest version}.
\item Family 291: \href{https://github.com/openai/math/blob/adc7f1241b42e322a6451854ab7e4b4c146bf78a/lean/docs/291.md}{Lean scope}.
\item \href{https://github.com/openai/math/blob/adc7f1241b42e322a6451854ab7e4b4c146bf78a/lean/ComparatorChallenges/UniformGamma.json}{Comparator UniformGamma.json}.
\item \href{https://github.com/openai/math/blob/adc7f1241b42e322a6451854ab7e4b4c146bf78a/lean/ComparatorChallenges/UniformGamma.lean}{Comparator UniformGamma.lean}.
\item \href{https://github.com/openai/math/blob/adc7f1241b42e322a6451854ab7e4b4c146bf78a/CONTENTS.md}{OpenAI catalogue entry}.
\end{itemize}
\end{document}
